% !TeX root = ../nc-quadric-moduli.tex
% 2026-11-22 / v0.1（非可換二次曲面の相対モジュライ原稿へ証明移管）
\subsection{二つの二元二次形式の不変式と good moduli}
\label{subsec:moduli-noncubic-invariants}

本節の $P$、$G$ は第~\ref{sec:moduli-boundary}節の二形式表示であり、
$\mathcal D=\mathbb P(\operatorname{Sym}^2V)\setminus V(\Delta)$ を
非退化二元二次形式の射影パラメータ空間とする。
二形式 $f=a u^2+buv+c v^2$、$g=a'u^2+b'uv+c'v^2$ に対し
\[
 \Delta(f)=b^2-4ac,\qquad D(f,g)=bb'-2ac'-2a'c,\qquad
 \chi(f,g)=\frac{D(f,g)^2}{\Delta(f)\Delta(g)}
\]
と置く。$\chi$ は二形式の基底変換・スカラー変換・交換に不変であり、
$P$ 上の正則関数である。
$\mathfrak N^{\mathrm{pres}}=[P/G]$ と書く。

\begin{proposition}[表示スタックの良モジュライ]
\label{prop:moduli-noncubic-good-moduli}
標数零の代数閉体 $k$ 上で
\[
 k[P]^G=k[\chi].
\]
従って $\chi:\mathfrak N^{\mathrm{pres}}\to\mathbb A^1_\chi$ は
Alper の意味の good moduli space である。
\end{proposition}
\begin{proof}
まず核の二次元トーラスの不変式を取ると、形式の代表の定数倍が除かれ、
\[
 k[P]^G
 =k[\mathcal D^2]^{\operatorname{PGL}_2\times S_2}
\]
となる。トーラスの作用が逆二乗であることは、標数零では不変式環を変えない。
$\operatorname{PGL}_2$ は $\mathcal D$ に推移的に作用するため、
第一形式を $uv$ と固定すれば、順序付き二形式の不変式環は
$uv$ の射影安定化群 $H=\mathbb G_m\rtimes C_2$ による
$\mathcal D$ の不変式環に同定される。
これは斉次空間 $\mathcal D=\operatorname{PGL}_2/H$ の faithfully flat な局所自明化で
不変関数を降下させることによる。

第二形式を $[a:b:c]$ と書く。
$\mathcal D$ は $\mathbb P^2$ の判別式非消滅開集合であり、アフィンで
\[
 k[\mathcal D]=\bigl(k[a,b,c,\Delta^{-1}]\bigr)_0,
 \qquad \Delta=b^2-4ac.
\]
残るトーラスの作用は射影座標で
$(a,b,c)\mapsto(t a,b,t^{-1}c)$ と書ける。
次数零のトーラス不変単項式は
$a^i b^j c^i\Delta^{-m}$、$2i+j=2m$ の形なので、$j$ は偶数である。
従って不変式環は
\[
 x=\frac{b^2}{\Delta},\qquad y=\frac{ac}{\Delta}
\]
で生成され、$x-4y=1$ により $k[x]$ である。
残る $C_2$ は $a,c$ を交換して $x$ を保つ。
第一形式が $uv$ なら $D(uv,f)=b$、$\Delta(uv)=1$ なので $x=\chi$ である。
二形式を交換する $S_2$ も対称な $\chi$ を保つため、全不変式環は $k[\chi]$ となる。

$P$ はアフィン、$G$ は標数零で線型簡約なので、商スタックから不変式環の
スペクトルへの射は good moduli space である
\cite[Example 8.3, Theorem 13.2]{Alper2009}。
\end{proof}

この計算は二形式の\emph{表示商}についての主張である。
表示商と内在的被約 noncubic 境界との相対同値、および
全内在的 $\mathfrak H_2$ における粗モジュライ不存在の証明は、
相対認識・自己同型の比較が必要なので別に扱う。
